%!TEX program = xelatex
% 如果你有参考文献（myref.bib），请按 xelatex -> bibtex -> xelatex*2 的顺序进行编译
\documentclass[dvipsnames, svgnames,a4paper,11pt]{article}
% Share-version redaction helpers
\newcommand{\answerblank}[1][3.8cm]{%
  \par\nopagebreak[4]\noindent\textbf{答：}\par\nopagebreak[4]\vspace*{#1}}
\newcommand{\recordblank}[1]{%
  \par\noindent\rule{0.95\linewidth}{0pt}\rule{0pt}{#1}\par}
\newcommand{\privateimage}[2]{%
  \rule{#1}{0pt}\rule{0pt}{#2}}
% ----------------------------------------------------
%   中山大学物理与天文学院本科实验报告模板
%   实验：全息照相实验
% ----------------------------------------------------

\input{settings} % 导入模板的相关设置


%---------------------------------------------------------------------
%	正文
%---------------------------------------------------------------------

\begin{document}

\scoresTable{30}{}{30}{}{20}{}{80}{} % 成绩栏（可由教师填写）

% \infoTable{专业}{年级}{姓名}{学号}{实验时间}{教师签名}
\infoTable{\rule{2.5cm}{0.4pt}}{\rule{2.5cm}{0.4pt}}{\rule{2.5cm}{0.4pt}}{\rule{3cm}{0.4pt}}{\rule{3cm}{0.4pt}}{\rule{3cm}{0.4pt}}

\begin{center}
	\LARGE 全息照相实验 \quad 实验报告
\end{center}

\textbf{【实验报告注意事项】}
\begin{enumerate}
	\item 实验报告由两部分组成：
	\begin{enumerate}
		\item 预习报告：课前认真研读\underline{\textbf{实验讲义}}，弄清实验原理；实验所需的仪器设备、用具及其使用，完成课前预习思考题；了解实验需要测量的物理量，并根据要求提前准备实验记录表格（可以参考实验报告模板，可以打印）。\textbf{（30分）}
	    \item 实验记录与分析：认真、客观记录实验条件、实验过程中的现象以及数据。实验记录请用珠笔或者钢笔书写并签名（\textcolor{red}{\textbf{用铅笔记录的被认为无效}}）。\textcolor{red}{\textbf{保持原始记录，包括写错删除部分，如因误记需要修改记录，必须按规范修改。}}（不得手记的值输入到电脑打印）；离开前请实验教师检查记录并签名。\textbf{（50分）}
	\end{enumerate}
	\item \textbf{本实验报告可提前打印出来，当场记录分析完成交给带实验的老师，课后无需再提交。}若当场完成不了，则请课后完成，再扫描并通过seelight提交。
	\item 实验中\textbf{避免激光器伤到眼睛}；避免用手直接接触镜片的光学面；安装镜片时需在光学平台上尽量靠近台面的高度操作；实验平台配件所用固定螺钉需拧紧，但不可过紧。
\end{enumerate}


\subsection*{实验安全注意事项}
\begin{enumerate}
	\item \textcolor{red}{\textbf{严禁激光直射眼睛}}，操作时须佩戴相应波长的防护眼镜，避免反射光入眼；
	\item 光学元件要轻拿轻放，\textbf{不得用手触摸光学面}，保持镜面清洁；
	\item 注意身体部位不要直接接触显影液、定影液等化学品，操作后及时洗手；
	\item 实验过程中尽量减少振动、噪声和气流干扰，\textbf{曝光期间绝对不能触碰光学平台};
	\item 全息干版对光极为敏感，在非必要情况下应保存在遮光盒内，装片操作须在全暗环境下进行。
\end{enumerate}

\clearpage


\setcounter{section}{0}
\nsection{全息照相}{全息照相实验}{预习报告}

\subsection{实验目的}
\begin{enumerate}
	\item 学习全息照相的基本原理和方法；
	\item 了解全息照相的主要特点；
	\item 学习全息照片的制作方法和技术；
	\item 学习观察全息照片的方法。
\end{enumerate}

\subsection{仪器用具}
\begin{table}[htbp]
	\centering
	\renewcommand\arraystretch{1.6}
	\begin{tabular}{p{0.05\textwidth}|p{0.20\textwidth}|p{0.05\textwidth}|p{0.55\textwidth}}
	\hline
	编号 & 仪器用具名称 & 数量 & 主要参数（型号，测量范围，测量精度等） \\
	\hline
	1 & 防震光学平台 & 1 &  \\
	2 & 氦氖激光器 & 1 & 波长 $\lambda = 632.8\,\text{nm}$，输出单横模（TEM$_{00}$）相干光，功率约数毫瓦 \\
	3 & 扩束透镜 & 1 & 短焦距凸透镜 \\
	4 & 分束器 & 1 &  \\
	5 & 反射镜 & 3 & 平面反射镜\\
	6 & 全息干版 & 1 & 高分辨率银盐感光干版，乳胶层厚度约 $6\sim15\,\mu\text{m}$ \\
	7 & 显影液 & 1 & D219 型显影液，使用温度 $20\,^\circ\text{C}$\\
	8 & 定影液 & 1 & 用于定影，处理时间 $2\sim4\,\text{min}$ \\
	9 & 暗房设备 & 1 & 暗绿灯、显影盘、定影盘、计时器等，供洗片使用 \\
	\hline
\end{tabular}
\end{table}

\subsection{原理概述}

\subsubsection{全息照相的基本思想}

全息照相核心在于利用光的干涉原理，将物光波前的\textbf{振幅}与\textbf{相位}信息同时完整记录在感光介质上，从而能够再现物体的真实三维图像。

与普通照相不同，普通照相以几何光学为基础，利用透镜在平面上记录光强分布，只保留了二维强度信息，丢失了相位信息，因此无法呈现真正的三维立体像。全息照相则借助参考光将物光的相位信息转化为感光介质可以记录的强度分布，从而实现对完整波前信息的记录和再现。

\subsubsection{全息图的记录原理}

激光经分束器后分为两路：一路照射被摄物体，经物体漫反射后到达全息干版，称为\textbf{物光}$a$；另一路直接（经反射镜导引）照射干版，称为\textbf{参考光}$r$。两束相干光在干版上相互叠加，产生干涉条纹。

设参考光在干版上某点的复振幅为 $r$，物光在该点的复振幅为 $a$，则干版上的光强分布为：
\begin{equation}
	I = |r + a|^2 = |r|^2 + |a|^2 + r^* a + r a^*
\end{equation}
其中，前两项为各自光强的直流分量，后两项为含有物光振幅与相位信息的干涉项。由于 $|r|^2$ 和 $|a|^2$ 在干版面上变化缓慢，记录介质主要感受的是随空间快速变化的干涉项，其结果是在干版上形成极细密的干涉条纹。条纹的间距、方向和对比度携带了物光的完整波前信息（振幅和相位），经曝光和化学处理后，这些信息被固化在感光乳胶中，形成全息图。

对于反射式全息照相，物光与参考光从干版两侧相对射入，夹角接近 $180°$，干涉极大形成的银层间距为：
\begin{equation}
	d = \frac{\lambda}{2\sin\theta} \approx \frac{\lambda}{2}
\end{equation}
以波长 $\lambda = 632.8\,\text{nm}$ 计算，银层间距约为 $0.3\,\mu\text{m}$，在乳胶层内可形成数十层结构，形成三维布拉格光栅。

\subsubsection{全息图的波前再现}

波前再现时，用与原参考光相同（或相近）的再照光照射已处理好的全息图。此时全息图相当于一个复杂的衍射光栅。记经曝光显影后干版的振幅透过率为：
\begin{equation}
	t = t_0 - \beta I = t_0 - \beta(|r|^2 + |a|^2) - \beta r^* a - \beta r a^*
\end{equation}
用再照光 $r$ 照射时，透射光中含有正比于 $r \cdot r^* a = |r|^2 a$ 的成分，即重建了原始物光波 $a$（相差一常数因子）。观察者从干版后方看去，便能看到与原物完全相同位置的\textbf{虚像}，具有完整的三维视差和深度感。

此外，透射光中还有正比于 $r \cdot r a^* = r^2 a^*$ 的共轭项，在一定条件下形成实像（孪生像），该实像位于全息图的另一侧，且有左右互换的特点。

\subsubsection{全息照相的特点}

\begin{enumerate}
	\item \textbf{三维性}：再现像具有完整的三维视差，改变观察角度可以看到物体不同侧面，并需重新调焦，与真实物体的视觉效果完全一致。
	\item \textbf{整体性}：全息图的任意局部小片均包含整个物体的信息，因为每点都受到物体各部分的物光和参考光的同时作用，故碎片也能再现完整图像，但分辨率随面积减小而降低。
	\item \textbf{多重记录性}：同一张干版可以通过多次曝光（改变参考光角度等）记录多幅独立的全息图，互不干扰，分别再现。
	\item \textbf{无正负像之分}：全息图复制品与原版再现的均为物体正像。
\end{enumerate}



\subsection{实验前思考题}

\begin{question}
	实验中应该注意哪些影响因素才能够保证成功观察到全息再现图像？
\end{question}
\answerblank[3.8cm]

\begin{question}
	物光和参考光的光程差应保持在什么范围？为什么？
\end{question}
\answerblank[3.8cm]

\clearpage
\begin{table}
	\renewcommand\arraystretch{1.7}
	\centering
	\begin{tabularx}{\textwidth}{|X|X|X|X|}
	\hline
	专业：& \rule{2.5cm}{0.4pt} & 年级：& \rule{2.5cm}{0.4pt} \\
	\hline
	姓名：&\rule{2.5cm}{0.4pt} & 学号：&\rule{3cm}{0.4pt} \\
	\hline
	室温：& \rule{2.5cm}{0.4pt} & 实验地点：& \rule{2.5cm}{0.4pt} \\
	\hline
	学生签名：&\rule{2.5cm}{0.4pt} & 评分：& \\
	\hline
	日期：& \rule{3cm}{0.4pt} & 教师签名：& \rule{3cm}{0.4pt} \\
	\hline
	\end{tabularx}
\end{table}

\nsection{全息照相}{全息照相实验}{实验记录}

\subsection{实验内容、步骤及结果}

\subsubsection{实验准备}
\begin{enumerate}
	\item 激光器开机预热约1小时，待输出稳定后开始搭建光路。
	\item 按仪器清单核对所有光学元件，检查各元件表面是否清洁完好。
	\item 在防震光学平台上依次安装激光器、分束器、扩束透镜、反射镜和干版架。
\end{enumerate}

\subsubsection{光路调整}

依照讲义对全息照相光路进行搭建：

\begin{enumerate}
	\item 激光器出射光经\textbf{分束器}分为两路；
	\item \textbf{参考光}：经反射镜导引，再通过扩束透镜扩束后，直接照射干版；
	\item \textbf{物光}：经扩束透镜扩束后照射被摄物体，物体漫反射后的散射光照射干版；
	\item 调整各元件位置，确保：
	\begin{itemize}
		\item 参考光与物光的夹角在 $30°\sim50°$ 之间；
		\item 两路光程大致相等，光程差控制在相干长度以内；
		\item 参考光与物光在干版处的光强比约为 $3:1\sim5:1$（用光电池配检流计在干版架位置测量）。
	\end{itemize}
	\item 用直尺测量各元件间距，按比例绘制实际光路图。
\end{enumerate}
\recordblank{5cm}

\subsection{实验光路图}

\recordblank{5cm}
\subsection{光路参数记录}
\recordblank{5.5cm}
\subsubsection{曝光与洗片}

\begin{enumerate}
	\item \textbf{装干版}：关闭室内所有光源，在全暗环境下，将药膜面朝向被摄物体方向固定于干版架。
	\item \textbf{曝光}：设定曝光10s时间进行曝光。曝光期间保持绝对安静，禁止走动和触碰台面。
	\item \textbf{显影}：在暗绿灯下，将干版放入D219显影液（温度约 $20\,^\circ$C）中，轻轻搅动，对暗绿灯观察至干版出现微黑时取出，用清水冲洗。
	\item \textbf{定影}：将显影后的干版放入定影液中定影，期间不断搅动定影液。
	\item \textbf{水洗}：定影完毕后，放入清水中冲洗。
	\item \textbf{干燥}：取出干版，自然晾干或用冷风吹干。
\end{enumerate}

\subsection{全息再现像照片}
\recordblank{6cm}
\subsection{原始数据记录}
\recordblank{5.5cm}
\subsubsection{图像再现观察}
\recordblank{5.5cm}
\subsection{实验过程中遇到的问题记录}
\recordblank{5cm}
\subsection{桌面整理}
\recordblank{5.5cm}
\clearpage
\subsection{实验后思考题}
\begin{question}
	与普通照相比较，全息照相有哪些特点？
\end{question}
\answerblank[3.8cm]

\begin{question}
	全息照相是如何把光波的相位记录下来的？
\end{question}
\answerblank[3.8cm]

\begin{question}
	实验中应该注意哪些影响因素才能够保证成功观察到全息再现图像？
\end{question}
\answerblank[3.8cm]

% ---------------------------------------------------------------------
%   参考文献
\clearpage
\phantomsection
\addcontentsline{toc}{section}{参考文献}
\bibliographystyle{unsrt}
\bibliography{myref}

\end{document}
